% TOP_PROBLEM: {"id": 30007048, "title": "Polynomiality of the natural invariant anticanonical index", "category_id": 5, "category": {"id": 5, "name": "algebraic_geometry", "display_name": "Algebraic Geometry", "description": "Geometric objects defined by polynomial equations.", "slug": "algebraic-geometry", "order_index": 5, "created_at": "2026-07-31T15:26:25.671Z"}, "status": "open"}
% FIRST_POSTED: 2026-10-06
% ENTRY_KIND: statement_only
% !TeX program = lualatex
\documentclass[11pt]{article}
\usepackage[margin=25mm]{geometry}
\usepackage{fontspec}
\setmainfont{Latin Modern Roman}
\usepackage{amsmath,amssymb,mathtools}
\usepackage{unicode-math}
\setmathfont{Latin Modern Math}
\usepackage{xurl,hyperref}
\hypersetup{colorlinks=true,urlcolor=blue}
\setcounter{secnumdepth}{0}
\setlength{\parindent}{0pt}
\setlength{\parskip}{5pt}
\setlength{\emergencystretch}{3em}
\newcommand{\R}{\mathbb R}
\newcommand{\N}{\mathbb N}
\newcommand{\Z}{\mathbb Z}
\newcommand{\Q}{\mathbb Q}
\newcommand{\C}{\mathbb C}
\newcommand{\D}{\mathbb D}
\newcommand{\F}{\mathbb F}
\newcommand{\sref}[2]{\hyperlink{source:#1:#2}{[S#2]}}
\newcommand{\cataloguescope}[1]{#1}
\begin{document}
\section*{Polynomiality of the natural invariant anticanonical index}\label{30007048}
\subsection{Definitions and mathematical statement}
\paragraph{Statement recorded in the supplied source.}
Let a compact real torus T act holomorphically on a compact complex manifold F and linearize $L=-K_F$ by the differential action. Put $I_T(m)=\Sigma _q(-1)^q$ dim $H^q(F,L^m)^T$. Must some positive a and rational polynomial P satisfy $I_T(m)=P(m)$ for every positive multiple of a, with $P(0)=I_T(0)$? No positivity or projectivity is assumed.
\subsection{Short English statement}
Let a compact real torus T act holomorphically on a compact complex manifold F and linearize $L=-K_F$ by the differential action. Put $I_T(m)=\Sigma _q(-1)^q$ dim $H^q(F,L^m)^T$. Must some positive a and rational polynomial P satisfy $I_T(m)=P(m)$ for every positive multiple of a, with $P(0)=I_T(0)$? No positivity or projectivity is assumed.
\subsection{Sources}
\begin{itemize}
\item[S1] OpenAI, manuscript citation page. Source of the formulation motivating this local problem.
\url{https://github.com/openai/math/blob/adc7f1241b42e322a6451854ab7e4b4c146bf78a/preprints/Invariant-anticanonical-indices-and-conversion-of-twisted-differentials-September-26-2026/README.md}.
\end{itemize}
\end{document}
